% Generated by scripts/assemble.py; corrections belong in source/corrections.json.
\examhead{87}{东南大学 2025 年数学分析试卷(回忆版)}{86}
% p087-b01; source PDF page 87
\par\Needspace{6\baselineskip}
\blockanchor{p087-b01}
\question{1}求空间曲线\par
\nopagebreak[4]
% p087-b02; source PDF page 87
\blockanchor{p087-b02}
\displayformula{\left\{ \begin{array}{l} x + y + z = 0 \dots \\ x ^ {2} + y ^ {2} = z ^ {2} \dots \end{array} \right.}
% p087-b03; source PDF page 87
\blockanchor{p087-b03}
的切线, 其满足平行于平面  $x + 2y = 0 \cdots$\par
% p087-b04; source PDF page 87
\blockanchor{p087-b04}
\question{2}$z = z(x, y)$  是由方程  $\cdots - \int_{y}^{-\sqrt{x^{2} + y^{2}}} e^{-t^{2}} dt - \cdots = 0$  确定，求  $\frac{\partial z}{\partial x}, \frac{\partial^{2}z}{\partial x \partial y}.$\par
% p087-b05; source PDF page 87
\blockanchor{p087-b05}
\question{3}求幂级数 $\sum_{n=1}^{\infty} (-1)^{n-1} n^2 (x - 1)^n$ 的收敛域与和函数.\par
% p087-b06; source PDF page 87
\blockanchor{p087-b06}
\question{4}求函数 $f(x) = \begin{cases} \frac{1}{x\ln 2} - \frac{1}{2^x - 1}, & x \neq 0; \\ \frac{1}{2}, & x = 0. \end{cases}$ 在 $x = 0$ 的导数.\par
% p087-b07; source PDF page 87
\par\Needspace{6\baselineskip}
\blockanchor{p087-b07}
\question{5}设 $f(x)$ 在 $[a, b]$ 上连续, 在 $(a, b)$ 上可导, 且 $f(x) \neq 0$, 证明: 存在 $\xi \in (a, b)$, 满足\par
\nopagebreak[4]
% p087-b08; source PDF page 87
\blockanchor{p087-b08}
\displayformula{\frac {f ^ {\prime} (\xi)}{f (\xi)} = \frac {1}{a - \xi} + \frac {1}{b - \xi} \cdot \cdot}
% p087-b09; source PDF page 87
\blockanchor{p087-b09}
\question{6}求 $\int_{L} \frac{x \, \mathrm{d}y - y \, \mathrm{d}x}{x^{2} + 4y^{2}}$, $L$ 是曲线 $(x - 1)^{2} + y^{2} = r^{2} (r \neq 1)$, 取正向.\par
% p087-b10; source PDF page 87
\blockanchor{p087-b10}
\question{7}求 $f(x, y, z) = x^3 + y^3 + z^3 - 2xyz$ 的最值, 其中 $x^2 + y^2 + z^2 = 4$.\par
% p087-b11; source PDF page 87
\blockanchor{p087-b11}
\question{8}求 $\iiint_{\Omega} |xy| z \, \mathrm{d}x \, \mathrm{d}y \, \mathrm{d}z$，其中 $\Omega$ 是由 $x + z = 1, x - z = 1, x + y = 1, x - y = 1, z = 0$ 围成的区域.\par
% p087-b12; source PDF page 87
\blockanchor{p087-b12}
\question{9}设 $F(x) = \int_{0}^{+\infty} e^{-y^2} \cos 2xy \, \mathrm{d}y$.\par
% p087-b13; source PDF page 87
\blockanchor{p087-b13}
(1) 证明: $2F'(x) + xF(x) = 0$.\par
% p087-b14; source PDF page 87
\blockanchor{p087-b14}
(2) 求 $F(x)$.\par
% p087-b15; source PDF page 87
\blockanchor{p087-b15}
\question{10}设 $f(x)$ 在 $[a, +\infty)$ 上一致连续, $g(x)$ 在 $[a, +\infty)$ 上连续, 且 $\lim_{x \to +\infty} [f(x) - g(x)] = 0$, 证明: $g(x)$ 在 $[a, +\infty)$ 上一致连续.\par
% p087-b16; source PDF page 87
\blockanchor{p087-b16}
\question{11}求 $\iint_{\Sigma} z \, \mathrm{d}S$, 其中 $\Sigma$ 是由 $x^{2} + z^{2} = 2x$ 与 $z = \sqrt{x^{2} + y^{2}}$ 为围成的有限曲面.\par
% p087-b17; source PDF page 87
\blockanchor{p087-b17}
\question{12}记 $u_{n}(x) = \frac{\ln(1 + n^{\alpha}x)}{n^{\beta}}, x > 0, \alpha, \beta > 0.$\par
% p087-b18; source PDF page 87
\blockanchor{p087-b18}
(1) 讨论 $\sum_{n=1}^{\infty} u_n(x)$ 的敛散性.\par
% p087-b19; source PDF page 87
\blockanchor{p087-b19}
(2) 在 $\sum_{n=1}^{\infty} u_n(x)$ 收敛的基础上, 讨论其一致收敛性.\par
% p087-b20; source PDF page 87
\blockanchor{p087-b20}
(3) 讨论 $\sum_{n=1}^{\infty} u_n(x)$ 的可导性.\par
% p087-b21; source PDF page 87
\par\Needspace{6\baselineskip}
\blockanchor{p087-b21}
\question{13}设 $f: \mathbb{R} \to \mathbb{R}$ 为连续函数, $f$ 无不动点, 记\par
\nopagebreak[4]
% p087-b22; source PDF page 87
\blockanchor{p087-b22}
\displayformula{f _ {n} (x) = \underbrace {f \circ f \circ f \circ \cdots \circ f} _ {n \text {个}} (x)}
% p087-b23; source PDF page 87
\blockanchor{p087-b23}
为 $f(x)$ 的 $n$ 次复合, 数列 $\{f_n(x)\}$ 是否有界?\par
